% supp_further.tex -- Supplementary Section S9: the second outbreak test and the analyses of contact records,
% tracer measurements and zone structure, in detail.
% Paths in "% src:" comments are relative to the repository root.  Every number is extracted from the stored
% results by code/v2_numbers.py (-> data/v2_numbers.json, v2_tab_*.tex); figures:
% code/v2_figures.py.
\section{The second outbreak test and the further analyses in detail}\label{appS_further:sec}

This section gives the details of the second outbreak test and of the analyses of contact records, tracer
measurements and zone structure (Section~\ref{M-s8_outbreaks:sec} of the main text). Their protocols, decision rules,
power and data are in Section~\ref{appS_prot2:sec}.

\subsection{Second outbreak test}\label{s8_outbreaks:sec:seconddetail}

\begin{figure}[tbp]
  \centering
  \includegraphics[width=0.82\textwidth]{v2_rr_forest_en}
  % generated by code/v2_figures.py (rr_forest) from data/bsc_validation2/a1_more_outbreaks/07_descriptive.json
  \caption{Near/far risk ratios of the eleven outbreaks of the second dataset, with 95\,\% intervals (a continuity
  correction of 0.5 where a stratum has no case), and random-effects pooled estimates. Near: within two rows
  of a source (aircraft), the index patient's bay or cubicle (ward), within 8\,m (processing line). Small
  symbols: ratio of the expected stratum attack rates under the constant-ratio model and under $\Mtwo$
  (Section~\ref{M-s8_outbreaks:sec:models}; for the six calibration events, from a fit to the other calibration
  events of the class).}
  \label{s8_outbreaks:fig:forest}
\end{figure}

\begin{figure}[tbp]
  \centering
  \includegraphics[width=\textwidth]{v2_holdout_strata_en}
  % generated by code/v2_figures.py (holdout_strata) from data/bsc_validation2/a1_more_outbreaks/10_corrected.json (primary_corrected.events)
  \caption{Second test: the five held-out SARS-CoV-2 events. Observed cases by distance stratum and the
  numbers expected, given the event total, under the well-mixed model, the constant-ratio model and $\Mtwo$,
  all calibrated on the six pre-2020 events. In brackets under each stratum: number of people at risk. On the
  four flights $\Mtwo$ and the constant ratio cannot be told apart; on the processing line the constant ratio
  uses a near zone of 8\,m, which the authors of the source selected on these same data.}
  \label{s8_outbreaks:fig:strata}
\end{figure}

Fig.~\ref{s8_outbreaks:fig:strata} shows the observed and expected cases by distance stratum for the five
held-out events. The two ward datasets contradict each other, and no model fits the inpatients in-sample.
% src: code/bsc_validation2/a1_more_outbreaks/PREREG_ADDENDUM_power.md (in-sample adequacy of W2: 0.001-0.025)
The room value is a correction made after the held-out events had been scored. The lattice propagator of the
first test, copied unchanged, loses precision at receivers whose exposure is below its round-off error,
which reaches about $10^{-10}$ of the largest exposure; the room posterior first computed (mode $0.02\sqmh$) came from two such noise values. A
propagator written as a positive sum of image terms removes the defect. The correction changes no number for
the flights, moves the pooled scores from $+13.37$ and $-10.13$ to the values reported below, and leaves the
outcome of the rule unchanged; the re-check confirmed it with a third propagator and found the
restaurant calibration of the first test unaffected (Section~\ref{appS_prot2:sec:numerics}). The hold-out
prediction for the classroom of the first test was affected and is reported corrected
(Section~\ref{M-s8_outbreaks:sec:holdout}).
% src: code/bsc_validation2/a1_more_outbreaks/POSTHOC_LOG.md (P5); data/v2_numbers.json (out2.cal.registered_room_M2: mode 0.020; out2.pooled_registered.M2.all: delta0 13.37, deltaC -10.13); notes/VERIFICATION.md (Section 4.6, second outbreak test, item 13: confirmed)

It is positive in each of the 15 pre-specified
sensitivity variants ($+10.3$ to $+18.1$). These variants were computed with the room kernel as first frozen,
before the numerical correction (primary values $+13.37$ and $-10.13$ with that kernel); the flight parts are
unaffected, and the pooled values would move by about $+1$.

It is significantly in favour of the constant
ratio in the three variants that change the outcome definition of the flights: $-2.52$ (mirror $p=0.013$) with
the confirmed and probable cases of the flight to Naha, $-1.90$ ($p=0.026$) with flight-associated cases only on
the flight from Sydney to Perth, and $-3.42$ ($p=0.005$) with all four changes of outcome definition. It is
significantly in favour of $\Mtwo$, $+1.21$ ($p=0.044$), when the members of two households seated together are
counted once. With the published ratio 2.4 \citep{hertzberg2016risk} in place of the calibrated one $\Mtwo$ is
ahead on the flights ($+1.18$, $p=0.028$) and behind over the five events ($-4.44$, mirror $p=0.002$, with
the room kernel as first frozen; $-4.21$, mirror $p=0.002$, with the re-check's propagator). Comparisons
added during the re-check give $+2.4$ ($p=0.003$) with two strata in place of four, and $+7.8$ and $+2.2$ with a
near zone of one and of five rows. On aircraft $\Mtwo$ is therefore not shown to be better than the two-row rule, nor shown
to be worse; the sign is not determined.

On the
processing line, $p=1.4\times10^{-5}$: the attack rate is flat out to 8\,m (5 of 9, 12 of 17) and then drops (1
of 22, 2 of 30), whereas the ward-calibrated kernel is almost well mixed and expects 2.9, 4.8, 5.7 and 6.5
cases. This is mostly a failure of the transfer from the ward to the line; at the coefficient that the line
itself prefers ($40\sqmh$) the log score of $\Mtwo$ is $-9.5$, still below the $-6.86$ of the constant ratio.

\emph{Secondary analyses of the protocol.} None changes the outcome, and two bear on what follows.
(i)~The kernel of the first test, calibrated on the train cohort ($\Dair\approx20\sqmh$), was applied to the
eight flights with no parameter fitted on any of them. It beats the well-mixed model ($\Delta_0=+7.56$) and is
rejected on four flights ($p=0.020$, 0.004, 0.0002 and 0.015): it is too steep for aircraft, although it
reproduces the two flights whose cases sat next to the sources ($p=0.72$ and 1.0).
(ii)~With the roles exchanged, calibrating on the SARS-CoV-2 events and predicting the earlier ones, $\Mtwo$
beats the constant ratio on the four pre-2020 flights ($\Delta_{\mathrm{C}}=+3.58$, $p=0.004$), because the
ratio learned on the 2020 flights, about 7, over-predicts the earlier gradients, above all on the influenza
flight with none; in the ward it loses ($-3.64$) and fails on the inpatients ($p=2\times10^{-6}$). Pooled,
$\Delta_{\mathrm{C}}=-0.06$.
(iii)~In the 15 sensitivity variants (outcome definitions, sets of sources, row pitch, cabin ventilation,
households, ward calibration) $\Delta_{\mathrm{C}}$ lies between $-12.8$ and $-7.7$ and the outcome is V5 in 12
and V3 (`better than well mixed, but a constant ratio predicts better') in three (kernel as first frozen, see
above).
(iv)~Leaving one event out within each class, $\Mtwo$ and the constant ratio are not separated on the
flights: $\Delta_{\mathrm{C}}=-0.06$, negligible in size but improbable under either model (one-sided
$p=0.046$ under the constant ratio, mirror $p=0.044$ under $\Mtwo$); formally the mirror criterion is met,
but the value is about as improbable under the constant ratio, and the re-check gives $+0.02$; in the rooms $\Mtwo$ is no better than the well-mixed model
($\Delta_0=-0.02$) and the constant ratio predicts better ($\Delta_{\mathrm{C}}=-14.0$).
(v)~Scored seat by seat on the three held-out flights with seat maps, the constant ratio does best ($-76.4$),
ahead of $\Mtwo$ ($-78.6$) and the well-mixed model ($-83.5$).
(vi)~The train kernel of the first test, applied to the rooms, is rejected on the ward inpatients
($p=2\times10^{-5}$) and on the processing line ($p=4\times10^{-4}$).
(vii)~The exploratory two-scale variant of the protocol ($\Mtwo$ plus a well-mixed fraction), evaluated by the
same rule, has $\Delta_0=+17.5$ and $\Delta_{\mathrm{C}}=-5.95$ (mirror $p=2\times10^{-4}$) and fails on the
processing line: outcome V3.
% src: data/v2_numbers.json (out2.sec.train_kernel.posterior: mode 19.95; out2.sec.train_kernel.flights: delta0 7.555; out2.sec.train_kernel.padeq: F1 0.0195, F4 0.0035, F5 0.0002, F6 0.0153, F7 0.72, F8 1.0; out2.sec.swap.air: delta0 6.42, deltaC 3.579, p_B2 0.0039; out2.sec.swap.room: deltaC -3.641; out2.sec.swap_W2: 1.6e-6; out2.sec.swap.all: deltaC -0.063; out2.check.swap.post.air|CRR: mode 7.08; out2.variants.summary: dC -12.78 to -7.65, n_V5 12, n_V3 3; out2.sec.loo_air: deltaC -0.058, p_B2 0.046, p_mirror 0.044; out2.check.loo_air: deltaC +0.018; out2.sec.loo_room: delta0 -0.025, deltaC -14.01; out2.sec.seat: CRR -76.40, M2 -78.62, M0 -83.49; out2.sec.train_kernel.rooms: W2 2.4e-5, P1 4.0e-4; out2.pooled.M3.all: delta0 17.55, deltaC -5.95, p_mirror 0.000215; out2.pooled.M3.decision: B3_failures [P1], outcome V3)

In the second test $\Mone$ was evaluated in its closed form only, as a further column of the same rule. Its
protocol fixed the reading in advance: the first test had shown with the exact simulator that the closed form
overstates the reach of $\Mone$ (on the bus its log score differs from that of the simulated model by 9.9 units,
$-9.58$ against $+0.30$ relative to $\Mzero$; pooled $-11.11$ against $-2.94$), so a pass of the closed form
would not rehabilitate $\Mone$ and a failure would confirm the first test. Against
the well-mixed model it passes with $\Delta_0=+15.45$ over the five held-out events, more than $\Mtwo$ ($+14.35$),
largely through the processing line ($+9.11$ against $+4.16$); against the constant ratio its pooled score is
$\Delta_{\mathrm{C}}=-8.06$ (mirror $p=1.5\times10^{-5}$), as for $\Mtwo$. It fails on three of the five
held-out events (flight to Naha $p=0.008$, flight EK448 $p=0.012$, processing line $p=0.004$), one more than
$\Mtwo$, and on the four flights the constant ratio predicts better ($\Delta_{\mathrm{C}}=-4.63$; a value this
low has probability 0.002 under $\Mone$), where $\Mtwo$ ties. Its formal outcome is the same as that of $\Mtwo$,
V5. Its gain over the well-mixed model is that of an approximation whose error on comparable seated events of
the first test is of the same size, and by the protocol it does not bear on the same-site model; its failures
are consistent with the first test. The same-site model is not supported by the first test, nor, as
Section~\ref{M-s8b_evidence:sec:contacts} shows, by recorded contacts.
% src: data/v2_numbers.json (out2.pooled.M1.all: delta0 15.446, p_B1 0, deltaC -8.055, p_mirror 1.5e-5; out2.pooled.M1.room: delta0 9.115 (out2.pooled.M2.room: 4.158); out2.hold.{F5,F7,P1}.M1.p: 0.0083, 0.0124, 0.0040; out2.pooled.M1.air: deltaC -4.634, p_mirror 0.0017; out2.pooled.M1.decision: B1 true, B2 false, B3_failures [F5, F7, P1], V5; out2.check.M1: air deltaC -4.62, p_mirror 0.0014); code/bsc_validation2/a1_more_outbreaks/PREREGISTRATION.md section 3 (M1cf); data/s8_first_test_exact.json (primary.events.T1: M1cf delta -9.585, M1 +0.297; primary.pooled: M1cf -11.109, M1 -2.939); data/bsc_validation/02b_m1_exact.json (C1 rows: D 0.021-14.93, median 0.2)

The second test confirms this with a test that its protocol lists. For the three room events a common
coefficient is rejected (likelihood ratio 17.4 on 2 degrees of freedom, approximate $p=1.6\times10^{-4}$): the students in
the ward prefer the lower end of the grid, the inpatients about $2{,}000\sqmh$ and the processing line
$40\sqmh$ (16--126). For the eight flights the heterogeneity is weaker (15.2 on 7 degrees of freedom,
approximate $p=0.034$; common value $158\sqmh$, 79--316). These $p$-values, too, use the chi-square reference
although one maximum lies on the edge of the grid and the events have 2 to 20 cases each; the per-event
profiles (Table~\ref{appS_prot2:tab:Dprofiles}) stand regardless. By pathogen the flights prefer $158\sqmh$ for SARS-CoV-2
(63--501), 126 for SARS-CoV-1 (40--794), 16 for tuberculosis (no lower limit) and about 500 for influenza,
where the gain over the largest coefficient of the grid ($10^{4}\sqmh$, close to well mixed) is 0.26 units
of log-likelihood, that is, no detectable gradient
(Table~\ref{appS_prot2:tab:Dprofiles}, Fig.~\ref{appS_prot2:fig:Dprofiles}). Together with the transfer of
the train kernel, which fails on four of eight flights (Section~\ref{M-s8_outbreaks:sec:second}), this leaves
no setting class with a single coefficient, and no transfer between classes.
Section~\ref{M-s8b_evidence:sec:tracer} compares the fitted coefficients with tracer fields, which correspond,
summarised by a diffusion coefficient, to values of order 100 to $1{,}000\sqmh$ in vehicles: far above the
vehicle value calibrated in the first test ($25\sqmh$); its room value, $0.5\sqmh$, is a few hundred times
below the 170 to $185\sqmh$ that a regression on two houses gives for rooms, for which no tracer measurement
exists. The values of the second test are not far below them: the aircraft value, 126, lies inside the range
measured in cabins (112 to 708), and the room value, 794, lies above the regression values. That does not make
an outbreak-fitted coefficient a measurement of air mixing, since within the same test no common coefficient
holds.
% src: data/v2_numbers.json (out2.D.room: LR 17.44, df 2, p 1.63e-4, per_event W1 0.01 (grid edge), W2 1995, P1 39.8 (15.8-125.9); out2.D.air: LR 15.18, df 7, p 0.0338, D_common 158.5, D_common_ci [79.4, 316.2]; by_pathogen: SARS-CoV-2 158.5 (63.1-501.2), SARS-CoV-1 125.9 (39.8-794.3), M. tuberculosis 15.8 (0.01-316), influenza 501 (50-1e4), gain 0.255, measured against the grid end 1e4 (code/bsc_validation2/a1_more_outbreaks/scripts/10_corrected.py: c.max() - c[-1]); out2.check.heterogeneity: room LR 17.67, air LR 15.00; out2.cal.air|M2 125.9; out2.cal.room|M2 794.3; tr.D.cabin: 112, 708; tr.D.table: classroom 169, office 185)

\subsection{Recorded contacts}\label{s8b_evidence:sec:contactsdetail}

\begin{table}[tbp]
  \centering\small
  \caption{Contacts generated by each model against RFID contact records in six indoor settings (calibrated on
  the first half of the days, scored on the second half). Families: number of the seven families of contact
  statistics passed, per dataset, in the order office 2013, primary school, hospital ward, office 2015, high
  school, conference (six families apply at the conference, which has no groups). Scenarios: SIR scenarios
  passed, of 24 (four per dataset). Ratio: secondary cases of the index case on model contacts divided by
  those on the recorded contacts, range over datasets of the mean over scenarios. Last column: mean absolute
  error of the attack rate over the 24 scenarios. The first four rows are the pre-specified test; the others
  are not.}
  \label{s8b_evidence:tab:contacts}
  % generated: data/v2_tab_contacts.tex (code/v2_numbers.py <- data/bsc_validation2/a2_contact_mobility/*.json, code/checks/contacts/check_summary.json)
  \begin{tabular}{@{}>{\raggedright\arraybackslash}p{0.34\textwidth}c>{\centering\arraybackslash}p{0.17\textwidth}ccc@{}}
    \toprule
    Model & Parameters & Families passed (of 7) & Scenarios & Ratio & Error \\
    \midrule
    \tabinput{data/v2_tab_contacts}
    \bottomrule
  \end{tabular}

  \smallskip
  \begin{minipage}{0.96\textwidth}\footnotesize
  $^{a}$\,Added post hoc during the re-check: contact rate of each pair taken from the training days, no space;
  exploratory.
  $^{b}$\,Developed on the first three datasets after the walk had failed on them; description after the fact
  (the confirmatory test on the last three datasets was uninformative).
  \end{minipage}
\end{table}

Two things are scored on the test days. \emph{Structure}: seven families of contact statistics (total contact
time; contact durations; gaps between contacts and recurrence of pairs; number of partners; heterogeneity
between people; repetition of pairs from one day to the next; share of contact within groups), each passed if
the model lies within a factor 0.8 to 1.25 of the data, or within a stated absolute margin. \emph{Epidemic
consequence}: an SIR epidemic is run on the recorded temporal network of the test days and on the
model-generated contacts, at the same total contact time, in four scenarios (nominal reproduction number 1.5
or 3; infectious period 1 or 4 days); a scenario is passed if the secondary cases of the index case agree
within 10\pct{} and the probability of a major outbreak and the mean attack rate within 0.05, and the epidemic
criterion is met with three of four scenarios. The rule: the movement component is validated on a dataset if
it meets the epidemic criterion and at least six of the seven families, and overall if it does so on
two-thirds of the datasets; the outcome is called A0 if it neither does this nor meets the epidemic criterion
where the well-mixed model fails. On synthetic data a true model passed in 20 of 20 replicates and a wrong one
reached the per-dataset criterion in none of 44. The protocol recorded the expectation A0.
% src: code/bsc_validation2/a2_contact_mobility/PREREGISTRATION.md sections 5-7; code/bsc_validation2/a2_contact_mobility/POSTHOC_LOG.md item 1 (0 of 44); data/v2_numbers.json (con.power.cells)

\begin{figure}[tbp]
  \centering
  \includegraphics[width=\textwidth]{v2_contacts_epidemic_en}
  % generated by code/v2_figures.py (contacts_epidemic) from data/bsc_validation2/a2_contact_mobility/{s1,s2b,s2d,s3a,s3b,s3d}_<dataset>.json (epi) and verify/v4_static_*.json
  \caption{SIR outcomes on model-generated contacts against the same outcomes on the recorded contacts of the
  test days: six datasets, four scenarios each. Open symbols: the three development datasets; filled: the
  three confirmatory ones. Shaded: tolerance of the protocol. Points above the diagonal are over-predictions.
  The lattice walk, in both forms, lies with the two models without space; the walks anchored to home sites
  were developed after that failure, and the static network was added during the re-check.}
  \label{s8b_evidence:fig:contacts}
\end{figure}

\paragraph{Result: the walk fails on all six datasets.}
The homogeneous walk passes one or two of the seven families on each dataset, namely those it was calibrated
on, the total contact time and the mean duration; the two-zone walk passes two on each
(Table~\ref{s8b_evidence:tab:contacts}; Fig.~\ref{appS_prot2:fig:structure}). What it does not reproduce is
who meets whom. People in the walk have too many partners on five of the six datasets: 6.0 per person and day
against 4.6 in the office of 2013, 34 against 15 in the high school, 90 against 47 in the primary school.
Contact within the recorded groups is 10 to 33\pct{} of contact time in the walk and 57 to 94\pct{} in the
data. Long contacts, the burstiness of the gaps between contacts (0.08--0.25 against 0.36--0.61) and the
repetition of the same pairs on the next day are missing as well. The real training days, used as a predictor
of the test days, reproduce the burstiness on all six datasets, the long contacts on four (not in the office of
2013 and the hospital ward) and the repetition of pairs on three of the four datasets where it is defined (not
in the hospital ward; it is undefined at the primary school and the conference).
% src: data/v2_numbers.json (con.model.RW0.S: [1,2,2,1,1,1]; con.model.RWhet.S: [2,2,2,2,2,2]; con.stats_obs_vs_RW0: S3_deg_mean 4.58/5.96 (InVS13), 14.6/33.8 (Thiers13), 46.6/90.4 (LyonSchool), 26.5/23.4 (SFHH); S6_within_frac observed 0.57-0.94, walk 0.10-0.33; S2_burst observed 0.36-0.61, walk 0.08-0.25; con.realtrain_S: [6,6,4,6,6,4]); data/bsc_validation2/a2_contact_mobility/s1_{InVS13,LyonSchool,LH10}.json and s3a_{InVS15,Thiers13,SFHH}.json (models.REALtrain.score: S2_intercontact true on all six; S1_duration false for InVS13, LH10; S5_persistence false for LH10, null for LyonSchool, SFHH)

The consequence for an epidemic is systematic. Because the walk spreads the same contact time over too many,
too equal partners, the index case infects more people on the walk's contacts than on the recorded ones in
all 24 scenarios, by 15 to 86\pct{}, and by 25 to 72\pct{} on average per setting; the mean attack rate is too
high by 0.10 to 0.29 (Fig.~\ref{s8b_evidence:fig:contacts}). No scenario is passed on any dataset. The total
contact time of model and data agrees within 2.5\pct{} in these runs, so the excess is not an artefact of the
contact rate. The re-check reproduced the result with its own data loader, walk simulator and
epidemic code, on files downloaded again from the source.
% src: data/v2_numbers.json (con.model.RW0: E_scen [0,0,0,0,0,0], n_cells_over 24 of 24, ratio_cell_min 1.151, ratio_cell_max 1.860, ratio_min 1.253, ratio_max 1.715, dA_per_dataset 0.103-0.292, contact_time_ratio 0.977-1.009); notes/VERIFICATION.md (Section 4.7, contact records, items 2, 3: confirmed)

The two models without space do neither better nor worse: both pass no scenario, with ratios of 1.25 to 1.79
(well mixed) and 1.20 to 1.63 (constant group ratio). The outcome is A0, with the sentence that the protocol
reserved for this case: the lattice walk adds nothing over a non-spatial contact model. The conclusion is, if
anything, understated. A static network with no space and no walk, in which each pair keeps the contact rate
it had on the training days, meets the epidemic criterion on three of the six datasets and passes 12 of the 24
scenarios, where neither pre-specified walk meets it on any dataset (the anchored walk with a slow own seat,
designed afterwards on the development datasets, meets it on one of them). This comparison was added during the
re-check after the results were known, with a crude estimator whose contact time is 2 to 9\pct{} low, so it is exploratory; and the
tolerance is tight, since the real training days used as a predictor of the test days meet the criterion on
four of the six datasets only. Neither point touches the result for the walk, whose error is several times
the tolerance.
% src: data/v2_numbers.json (con.model.WM: ratio 1.248-1.793, scen_total 0; con.model.BLOCK: 1.201-1.631, 0; con.ladder: A0 for every model; con.check.static: scenarios 0,4,0,3,2,3 -> 12 of 24, 3 of 6 datasets; con.model.RW0.E_datasets 0, con.model.RWhet.E_datasets 0, con.model.RWstickO.E_scen [0,0,3,2,0,0], E_datasets 1; con.check.train_replay: 2,4,4,3,1,4 -> 4 of 6 datasets); notes/VERIFICATION.md (Section 4.7, contact records, items 6 and 9)

One further property is against the mass-action picture of the model. In the walk, as in the well-mixed
model, the number of contacts grows as the square of the number of people present (fitted exponents 1.95 to
2.20). In the data the exponent is 0.45, 0.95, 0.76 and 0.81 in the two offices, the high school and the
conference, and 2.8 in the hospital ward: in four of five settings a person's contact rate does not grow
with occupancy. This comparison is descriptive: occupancy varies with the time of day, so the exponent
mixes density with schedule, and 20-minute intervals without a contact were left out of the fit, which biases
the observed exponent on the sparse datasets.
% src: data/v2_tab_datasets.tex (S7_alpha observed / walk: 0.45/1.97, 2.83/2.20, 0.95/2.06, 0.76/2.12, 0.81/1.95; primary school not comparable, occupancy range 1.27 < 1.5); notes/VERIFICATION.md (Section 4.7, contact records, item 7: partially confirmed; further point 9: bins with zero contacts dropped)

\paragraph{Walkers anchored to home sites: neither supported nor refuted.}
The protocol provided for a second stage. On the three development datasets the smallest changes that repair
the failed families were sought, keeping a Markov walk with same-site contact: a home site for each person
with a drift towards it, homes of the same group clustered in space, and either a slow seating zone, which is
the heterogeneous mobility $D_x$ of the model, or a walker that is slow on its own seat only, which is not.
The three variants, with their automatic fitting procedure, were frozen before the three confirmatory datasets
were touched. On those datasets none meets the epidemic criterion (the variant with a slow own seat passes
2, 0 and 0 of 4 scenarios; the variant with a slow seating zone passes none of 12), and the formal outcome is
again A0.
% src: code/bsc_validation2/a2_contact_mobility/POSTHOC_LOG.md items 3-13; code/bsc_validation2/a2_contact_mobility/FIX_FREEZE.json; data/v2_numbers.json (con.model.RWstickO.E_scen: [0,0,3,2,0,0]; con.model.RWstickZ.scen_conf 0; con.model.RWclus.scen_conf 0)

This outcome must not be read as a test that anchoring failed. In the re-check synthetic
contact data were generated from the anchored models themselves and the frozen fitting procedure was applied to
them: it failed
the epidemic criterion in six of seven replicates, with degenerate fits in four. A procedure that cannot
recover its own truth cannot refute a model, so the confirmatory stage is uninformative. What remains is a
description made after the fact. Over all 24 scenarios the mean absolute error of the attack rate falls from
0.17 for the homogeneous walk to 0.04--0.05 for the anchored variants, and the error of the index case's
secondary cases by 47 to 70\pct{}; two of the three variants still over-predict the secondary cases in all 24
scenarios, by 6 to 36\pct{}. These figures use the recorded group labels and are not at matched contact
time: the ratio of model to recorded contact time is 0.96 to 1.05 and 0.92 to 1.06 for the two variants with
clustered homes, and 0.98 to 1.11 for the variant with a slow own seat except in the high school, where it is
2.34; the reduction of 70\pct{} belongs to that variant and includes the high school. They indicate a
direction, home anchoring with clustered homes, and are not a result; no variant reproduces burstiness or the
repetition of pairs between days. A random walk in which people slow down near attractive individuals was
proposed and compared with SocioPatterns records by \citet{starnini2013modeling}; the anchored walks found here
by trial are close to that idea, and what is added is the epidemic consequence and the comparison with models
without space.
% src: notes/VERIFICATION.md (Section 4.7, contact records, items 4, 5; further points 1, 4, 5); data/v2_numbers.json (con.check.self_power; con.model.RW0.mean_abs_dA 0.167, mean_abs_log_ratio 0.361; con.model.RWclus: 0.041, 0.191, ratio_cell 1.07-1.36, n_cells_over 24; con.model.RWstickZ: 0.038, 0.166, 1.06-1.35, 24; con.model.RWstickO: 0.050, 0.108, n_cells_over 17; contact_time_ratio: RWclus 0.960-1.052, RWstickZ 0.917-1.058, RWstickO 0.983-1.109 and 2.340 (Thiers13); 1 - 0.108/0.361 = 0.70)

\paragraph{Scope.}
The result holds for one identification of the walk: a site is a contact range, steps are synchronous, and the
lattice has as many sites as the contact rate requires. The calibrated hop probabilities, 0.20 to 0.42 per
20-second step, correspond to a mobility of about 220 to 450 sites$^2$ per day if a site is identified with a
cell of Section~\ref{M-s2_model:sec}: between $\Dz=1\cellday$, which was not tested against contacts, and walking
mobility. Other site sizes and the continuous-time walk were not tried. The size of the over-prediction depends on the design with a nominal reproduction number, which
forces a high transmission probability per contact interval on sparse datasets; it is smallest in the primary
school (ratios 1.15 to 1.34). The sensors record proximity, not position, so the occupancy of zones could not
be tested, and neither could a measured mobility field.
% src: notes/VERIFICATION.md (Section 4.7, contact records, further points 6, 7); code/bsc_validation2/a2_contact_mobility/POSTHOC_LOG.md item 17; data/v2_numbers.json (con.dataset.*.p: 0.203-0.417; per-direction hop rate (p/4) x 4,320 steps per day = 219-451; data/bsc_validation2/a2_contact_mobility/s1_LyonSchool.json (epi.*.RW0.R_index / epi.*.REAL.R_index: 1.151, 1.321, 1.204, 1.337, ratios computed from the stored values)); code/bsc_validation2/a2_contact_mobility/PREREGISTRATION.md section 3 (move to a uniformly chosen neighbour accepted with probability p per 20 s step)

\subsection{Tracer measurements}\label{s8b_evidence:sec:tracerdetail}

\begin{table}[tbp]
  \centering\footnotesize
  \setlength{\tabcolsep}{3.4pt}
  \caption{The kernel fixed by tracer measurements, with nothing fitted to outbreaks except one
  $\varepsilon_k$ per event, on the stratum events of the first dataset. RR: predicted and observed near/far risk
  ratio; $p$: two-sided predictive $p$ of the observed near count under the tracer kernel and under
  $\Mzero$; last three columns: log score of the tracer kernel minus that of $\Mzero$ and of the constant
  ratios 2 and 3.}
  \label{s8b_evidence:tab:tracer}
  % generated: data/v2_tab_tracer.tex (code/v2_numbers.py <- data/bsc_validation2/a3_tracer_physics/04_evaluation.json)
  \begin{tabular}{@{}lccccccc@{}}
    \toprule
    Event & Near vs far & RR pred.\ / obs. & $p$ & $p$, $\Mzero$ & $-\Mzero$ & $-$ratio 2 & $-$ratio 3 \\
    \midrule
    \tabinput{data/v2_tab_tracer}
    \midrule
    Pooled & & & & & $+7.87$ & $-0.80$ & $-0.78$ \\
    \bottomrule
  \end{tabular}

  \smallskip
  \begin{minipage}{0.96\textwidth}\footnotesize
  $^{a}$\,With a continuity correction (no case at the remote tables). $K$: number of cases counted as
  acquired at the two neighbouring tables; with the source's own count, $K=5$, the tracer kernel is rejected
  ($p=0.002$).
  \end{minipage}
  % src: data/v2_numbers.json (tr.pooled_sum: M0 +7.873, CRR2 -0.800, CRR3 -0.780; tr.restaurant_byK: P 0.095, 0.027, 0.0074, 0.0019 for K = 2..5)
\end{table}

\paragraph{Measured mixing coefficients.}
If the non-local term of $\Mtwo$ describes the transport of aerosol, its mixing coefficient can be measured
without any outbreak, by releasing a tracer and recording its concentration. Published measurements exist for
four of the settings of Section~\ref{M-s8_outbreaks:sec}: ethane tracer gas at eight seats of the Hunan coach
itself, with a tracer-validated flow simulation for all seats \citep{ou2022}; fluorescent 1\,\textmu m particles
released from a breathing mannequin in the forward cabin of a wide-body aircraft \citep{kinahan2021aerosol};
salt aerosol along the saloon of an inter-city rail carriage \citep{woodward2022evaluation}; and tracer gas at
seven tables of the Guangzhou restaurant, with the source's flow simulation at the others \citep{li2021}. For the classroom and the open-plan office no
measurement exists, and the coefficient was taken from a regression of the turbulent diffusion coefficient on
the air change rate measured in two houses \citep{cheng2011modeling}, whose measured values are 3.6 to
$47\sqmh$ at below 0.2 to above 5 air changes per hour; it is extrapolated to rooms of several hundred cubic
metres at 2 to 12 air changes per hour. The estimators were fixed
before any coefficient was computed: a least-squares fit of the logarithm of the steady concentration of
$\Mtwo$ to the logarithm of the tracer reading, on the lattice of the event.
% src: code/bsc_validation2/a3_tracer_physics/PREREGISTRATION.md section 3; data/tracer/SOURCES.md (Cheng et al. 2011: K from 0.001 to 0.013 m2/s, i.e. 3.6 to 46.8 m2/h); data/v2_numbers.json (tr.D.table: classroom 429 m3, 2-12 ACH, 169 (38-745))

\begin{figure}[tbp]
  \centering
  \includegraphics[width=\textwidth]{v2_tracer_kernels_en}
  % generated by code/v2_figures.py (tracer_kernels) from data/tracer/{ou2022_B1_seats.csv, kinahan_tidy.csv, woodward2022_fig9_points.csv}
  % and data/bsc_validation2/a3_tracer_physics/{01_kernels.json, 05_descriptive.json}; data/bsc_validation2/a1_more_outbreaks/10_corrected.json (primary_corrected.post.air|M2)
  \caption{Tracer measurements. (a)~Coach of the Hunan outbreak: ethane measured at eight seats and the
  tracer-validated flow simulation at all seats \citep{ou2022}. (b)~Forward cabin of a wide-body aircraft:
  time-integrated particle counts by row for two releases, zero readings not shown
  \citep{kinahan2021aerosol}. (c)~Rail carriage: salt aerosol along the saloon, release at the vertical line
  \citep{woodward2022evaluation}; points digitised from the published figure. The lines in (a) and (c) are
  exponentials with the decay length $\sqrt{\Dair/\remrate}$ of the measured coefficient and of the $25\sqmh$
  fitted to the train cohort; they illustrate the two length scales and are not the fitted lattice fields.
  (d)~Mixing coefficients: measured (with range), fitted to outbreaks in the first test, fitted to the four
  pre-2020 flights of the second test (with 90\,\% interval), and preferred by the outbreak of that setting
  taken alone.}
  \label{s8b_evidence:fig:kernels}
\end{figure}

In the three vehicles the measured fields, summarised by a diffusion coefficient, are of order 100 to
$1{,}000\sqmh$ (Fig.~\ref{s8b_evidence:fig:kernels}d): $1{,}259$ in the coach (5th percentile of a bootstrap 708;
the eight points are nearly flat and bound the coefficient from below only), 708 and 112 for the two releases
in the cabin (geometric mean 282), and 891 and 355 for a release in the middle and at the end of the carriage.
For the classroom and the office no tracer measurement exists, and no coefficient was derived from the
restaurant's readings, which enter as a kernel directly; the regression extrapolated to the volume and
ventilation of the classroom and the office gives about 170 to $185\sqmh$ (38 to 753). The coefficients fitted to outbreaks in the
first test are $25\sqmh$ for vehicles and $0.5\sqmh$ for rooms, and the flight VN54 alone prefers 0.06 to 0.1.
The measured vehicle values are thus 4.5 to about 50 times the vehicle value of the first test, and the
extrapolated room values a few hundred times its room value. The second test is different: its aircraft
coefficient, $126\sqmh$ (52--438), lies inside the range measured in cabins, and its room coefficient, $794\sqmh$
(302--4{,}390), lies above the extrapolated room values. Three reservations on the cabin value: the fit uses positive
readings only, and estimators that keep the zero readings give 56 to $450\sqmh$; the two fits move from 708
and 112 to 891 and $178\sqmh$ if four additional sensors, one of them with an outlying reading, are left out;
and the measured field is not diffusive (it is biased forwards and cut off aft of row 8), so that a single
coefficient summarises it loosely.
The decay lengths $\sqrt{\Dair/\remrate}$ are about 2 to 5\,m in the cabin and 5 to 14\,m in the carriage and
the coach.
% src: data/v2_numbers.json (tr.D.coach: D 1259, lo 708; tr.D.cabin: D_5L 708, D_5A 112, D_geomean 282; tr.D.train: middle 891, end 355, len 8.3, 5.2; tr.D.table: classroom 169, office 185, coach length 13.9, cabin length 2.8; out2.cal.air|M2: 126, 52-438; out2.cal.room|M2: 794, 302-4391; tr.D.table: classroom 169 (38.4-745.2), office 185 (45.6-753.0)); data/checks/tracer.json: cabin_D (56-450 with zeros kept; 112 -> 178 and 708 -> 891 without the four additional sensors), cabin_decay_lengths_m (1.8-4.5), ratio_tracer_to_outbreak_D_at_low_cabin_value (4.5); notes/VERIFICATION.md (Section 4.8, tracer measurements, item 12; further point 3: field not diffusive; further point 4)

\begin{figure}[tbp]
  \centering
  \includegraphics[width=\textwidth]{v2_tracer_predictions_en}
  % generated by code/v2_figures.py (tracer_predictions) from data/bsc_validation2/a3_tracer_physics/04_evaluation.json (events.*.models.*: mean_near, int95)
  \caption{Number of cases in the near stratum, given the event total: mean and 95\,\% predictive interval
  under the well-mixed model, the kernel fixed by tracer measurements, the constant ratios 2 and 3, and
  $\Mtwo$ with the coefficient fitted to outbreaks in the first test (at its posterior mode; the restaurant was
  a calibration event). The tracer kernel removes the over-prediction on the coach by coming close to the
  well-mixed model. The gradients of flight VN54 and of the restaurant are not reproduced: the formal
  non-rejection of the flight depends on zero sensor readings (see text), and the restaurant is shown at the
  lower bound of its case count.}
  \label{s8b_evidence:fig:tracerpred}
\end{figure}

\paragraph{Confronted with the outbreaks.}
With the kernel fixed in this way, the seven events of the first dataset were predicted with one free parameter
each, the infectiousness (Table~\ref{s8b_evidence:tab:tracer}, Fig.~\ref{s8b_evidence:fig:tracerpred}). The
trains and the restaurant, calibration records before, are test events here.

\emph{Coach and call centre: no longer rejected, and not informative.} For the coach the measured kernel
predicts a rear/front ratio of 1.29 where 1.03 was observed ($p=1.00$); the kernel fitted to outbreaks had
predicted 11. For the call centre it predicts little desk-scale clustering, and the probability of a join-count
$z$ as small as observed is 0.10 (Fig.~\ref{s8_outbreaks:fig:callcentre}c). In both cases the measured kernel
is close to the well-mixed model, which fits as well or better ($p=1.00$ and 0.26); for the call centre the
kernel still predicts clustering of the wrong sign (median $z=+0.78$ against $-0.64$ observed). The two
failures of Section~\ref{M-s8_outbreaks:sec:holdout} are removed only in the sense that the model has stopped
predicting a gradient.
% src: data/v2_numbers.json (tr.event.T2: P p 1.0, rr_expected 1.285, M0 p 1.0, M2cal_mode p 0.0029; tr.callcentre: P P_le_obs_pooled 0.101, median 0.785; M0 0.258; z_obs -0.645); notes/VERIFICATION.md (Section 4.8, tracer measurements, items 4, 5)

\emph{Bus: inconclusive.} Predicted ratio 1.03 against 1.60 observed, $p=0.22$; the well-mixed model has
$p=0.20$.
% src: data/v2_numbers.json (tr.event.T1: P p 0.217, rr_expected 1.033; M0 p 0.204)

\emph{Classroom: the one gain.} The room relation predicts a front/back ratio of 1.77 where 2.8 was observed
($p=0.42$), and the well-mixed model is rejected ($p=0.036$). The result holds when the coefficient is doubled
or halved and under two other dose--response laws, and the implied emission, 317 quanta per hour (173--642),
lies within the range of the single events of Table~\ref{s8_outbreaks:tab:levels}, unlike the median
$8\times10^{7}$ that the first test needed. Constant ratios of 2 and 3 do as well, and the coefficient is
extrapolated from two houses, with the teacher's position assumed.
% src: data/v2_numbers.json (tr.event.C1: P p 0.424, rr_expected 1.769, M0 p 0.036, KB_Dx2 0.223, KB_Dhalf 1.0, P_lin 1.0, P_gam 0.228, CRR2 0.663, CRR3 1.0; tr.emission."C1 classroom": [172.9, 317.5, 642.2])

\emph{Flight VN54: not reproduced.} The predicted ratio is 1.66 against 7.3 observed: with 11 of 12 near
passengers infected and 1 of 8 far, the exponential dose--response law needs a dose ratio of about 19, and the
measured kernels give 1.5 to 3.9. The formal test does not reject ($p=0.14$), but only because the predictive
distributions of two releases that contradict each other are averaged, and because three sensors that map onto
far seats read exactly zero in one of them. With the zeros treated as missing, a variant listed in the
protocol, $p=0.003$; the flight is rejected in four of the six listed variants. The tracer tests were made in
another airframe.
% src: data/v2_numbers.json (tr.event.T5: P p 0.136, rr_expected 1.660, KA_5L 0.0007, KA_5A-mirrored 0.229, KA_zeros_missing 0.0028, KB_measD 0.045, KB_measD_oldvent 0.029, P_gam 0.0496, P_lin 0.185; tr.dose_ratio.T5: 18.6); data/bsc_validation2/a3_tracer_physics/01_kernels.json (cabin.*.direct.*.near_far_mean_ratio: 1.49-3.86); code/bsc_validation2/a3_tracer_physics/POSTHOC_LOG.md (P4); notes/VERIFICATION.md (Section 4.8, tracer measurements, item 8: to be reported as not reproduced)

\emph{Restaurant: not reproduced at the reported count.} The tracer concentration at the 13 remote tables
was on average 0.55 of that at the index table (0.32 to 0.86); four of these values are measured and nine are
the source's flow simulation, the lowest among them. With two cases counted as acquired at the neighbouring tables
the kernel is not rejected ($p=0.095$); with the five that the source reports, against none among 63 remote
patrons, it is ($p=0.002$).
% src: data/v2_numbers.json (tr.restaurant_tracer: remote_mean 0.553, 0.32-0.86; tr.restaurant_byK: 0.095, 0.027, 0.0074, 0.0019); data/bsc_outbreaks/tables/li2021_restaurant_tables.csv (norm_measured_tracer at TB, TC, T10, T15, T16, T17, T18; norm_predicted_tracer used elsewhere, code/bsc_validation2/a3_tracer_physics/src/a3lib.py r1_tables; remote minimum 0.32 at T09, simulated)

\emph{Trains: failure.} The kernel measured in the carriage is nearly flat over three rows and expects 14.6
cases in the row of the index case where 50 were observed (deviance 112.8 on 21 degrees of freedom,
$p<0.0005$). Every model fails on this table: the well-mixed model (125.6), constant ratios of 2 and 3 for the
index row (81.5 and 64.1) and the kernel that was fitted to the table in the first test (72.3).
% src: data/v2_numbers.json (tr.trains: P deviance 112.83, M0 125.61, CRR2 81.51, CRR3 64.06, M2cal_mode 72.30; tr.trains_rows.0: k 50, P 14.56)

\emph{Pooled.} Over the five stratum events the tracer kernel scores $+7.87$ against the well-mixed model
(exact one-sided $p=1.4\times10^{-6}$; $+3.41$ without the flight) and $-0.80$ and $-0.78$ against the constant
ratios 2 and 3, which themselves beat the well-mixed model by more ($+8.67$ and $+8.65$). The outcome of the
rule is S4, one adequacy failure (the trains) and no gain over the generic gradient; on the defensible reading
of the flight it is S5, two failures. The formal outcome S4 does not discriminate between the truths considered
in the power analysis (Table~\ref{appS_prot2:tab:tracerrule}, Section~\ref{appS_prot2:sec:tracer}). A
comparison added during the re-check, after the results were known, gives the probability of the pooled pattern
alone (better than well
mixed, not better than both constant ratios): 0.23 if the tracer kernel is true, 0.84 under a constant ratio of
2 and about 1.00 under a ratio of 3, so that the pooled evidence favours a generic ratio over the kernel by a
factor of about four (3.7 to 4.4).
% src: data/v2_numbers.json (tr.eval.pooled.P: M0 delta 7.873, p 1.44e-6; CRR2 -0.800; CRR3 -0.780; tr.crr_vs_M0: 8.673, 8.653; tr.eval.decision: outcome S4, D1_failures [T4]; tr.pooled_without.pooled_P_vs_M0_without_T5: +3.411); data/v2_numbers.json (tr.power.pattern: P_D2_not_D3 P 0.229, CRR2 0.842, CRR3 0.997; ratios 3.68, 4.36; enumerated by code/v2_numbers.py from data/bsc_validation2/a3_tracer_physics/02_predictions.json); data/checks/tracer.json: pooled_pattern_probability (0.23 / 0.84 / 1.00); notes/VERIFICATION.md (Section 4.8, tracer measurements, item 11: outcome S5 if the flight is counted as a failure)

\paragraph{An exploratory extension to four flights.}
After these results were known, the same cabin coefficients (112 and $708\sqmh$, equal weights) were applied
to the four held-out flights of Section~\ref{M-s8_outbreaks:sec:second}. The extension has a protocol of its
own, written with the near/far counts and the outcome of the second outbreak test known, and it reverses two
exclusions of the main protocol; it is exploratory and takes no part in any decision. The tracer kernel is not
rejected on any of the four flights ($p=0.054$, 0.26, 0.25 and 1.00), scores $+10.4$ against the well-mixed
model and $+2.8$ against a constant ratio of 2 ($p=0.005$), and ties with a ratio of 3 ($-0.16$). The absence
of a rejection is at the edge: with the geometric-mean coefficient that the main protocol specifies, the
flight to Naha is rejected ($p=0.044$), and each of the two coefficients taken alone fails one flight
($p=0.025$ and 0.002). The tracer measurements were made in wide-body cabins, and two of the four flights
were on Boeing~737 aircraft (Table~\ref{appB_dataset:tab:second}), so the transfer of the coefficient is
itself untested. What survives, as an exploratory result, is small: the two cabin coefficients (112 and
$708\sqmh$), measured without any outbreak, beat the well-mixed model and tie with a constant ratio of 3.
% src: data/v2_numbers.json (tr.ext.events.*.models.P.p: 0.054, 0.260, 0.246, 1.0; P_D708 on F7: 0.025; P_D112 on F5: 0.0016; tr.ext.pooled.P: M0 +10.36, CRR2 +2.83 (p 0.005), CRR3 -0.16); code/bsc_validation2/a3_tracer_physics/PREREG_EXT_flights.md section 0; data/checks/tracer.json: extension_naha.p_with_single_geometric_mean_D (p = 0.044, D = 282); notes/VERIFICATION.md (Section 4.8, tracer measurements, item 13; further point 9: tracer data are wide-body only); data/v2_tab_events2.tex (F5 B737-800, F8 B737-900)

\paragraph{Reading.}
Measured mixing is fast. A kernel with the measured coefficient is therefore close to well mixed in the coach
and the carriage and mild in the rooms, and it cannot produce the steep gradients of the train cohort, of
flight VN54 or of the restaurant. Those gradients need something other than diffusion at the scale of the
enclosure: exposure within about a metre, companions, directed airflow. The tracer kernel is not validated as
a spatial predictor; its one gain over the well-mixed model, the classroom, is matched by a constant ratio.
Only the coach and the restaurant have a tracer test in the enclosure of the outbreak; the cabin, the carriage
and the rooms rest on transfers that are themselves untested.
% src: notes/VERIFICATION.md (Section 4.8, tracer measurements, summary; further point 9)

\subsection{Zone-level structure}\label{s8b_evidence:sec:zonesdetail}

\begin{table}[tbp]
  \centering\small
  \caption{Zone-level test: influenza in 29 primary schools, with mixing shares measured in another school.
  Differences are of held-out log-likelihood (two-fold cross-validation by school), with 95\,\% bootstrap
  intervals over schools. The last two rows are post hoc, added during the re-check.}
  \label{s8b_evidence:tab:zone}
  % src: data/v2_numbers.json (zone.D_Z_minus_H: 798.4 [561.9, 1072.5]; zone.D_Z_minus_X: 211.1 [109.8, 340.4]; zone.D_Z_minus_C: 18.8 [-8.8, 54.0]; zone.primary_other: w_hat [0.659, 0.104, 0.236], LR_F_vs_Z 42.71, TV 0.1028, disp_obs 3.383, disp_pred_Z [2.65, 3.17, 3.61]; zone.TV_boot_ci [0.047, 0.168]; zone.check.hazard_bound."1e-09".D_N702: -19.1 [-37.8, -3.3]; zone.check.D_Z_minus_H: 744.8 [530.7, 999.0])
  \begin{tabular}{@{}>{\raggedright\arraybackslash}p{0.47\textwidth}>{\raggedright\arraybackslash}p{0.27\textwidth}>{\raggedright\arraybackslash}p{0.19\textwidth}@{}}
    \toprule
    Criterion & Result & Outcome \\
    \midrule
    Tier 1: zone model minus well-mixed school & $+798$ (562 to 1{,}073) & met \\
    Tier 1: zone model minus independent classes & $+211$ (110 to 340) & met \\
    Shares own class / grade / school: measured; fitted & 0.76 / 0.10 / 0.13; 0.66 / 0.10 / 0.24 & \\
    Tier 2a: equality of the shares (likelihood ratio, limit 5.99) & 42.7 & not met \\
    Tier 2b: total-variation distance (limit 0.15) & 0.103 (0.047 to 0.168) & met, not securely \\
    Tier 3: between-class heterogeneity, observed and interval of the zone model & 3.38; 2.65 to 3.61 &
      met; independent classes also contain it \\
    \addlinespace
    Zone model minus two-level model with one fitted ratio & $+18.8$ ($-8.8$ to $+54.0$) & not distinguishable \\
    Zone model minus zone model with guessed shares 0.7 / 0.1 / 0.2 & $-19$ ($-38$ to $-3$) & the guess is better \\
    Tier 1 without the lower bound on the external hazard & $+745$ (531 to 999) & unchanged \\
    \bottomrule
  \end{tabular}
\end{table}

The zone model takes these three shares as given and fits the overall rate $\beta$ and two parameters of an
external hazard from the city, in a daily chain-binomial model whose transmission kernel is the gamma
density with mean 1.7 and standard deviation 1.0 days used by \citet{endo2021within}, chosen by them to give
the mean serial interval of 2.2 days of influenza A(H3N2) \citep{vink2014serial}. Its competitors are fitted in the same way: a well-mixed school, independent classes,
and a generic two-level model with one fitted ratio of own-class to rest-of-school risk, the analogue of the
constant-ratio model. A model with three free shares serves as a ceiling. The rule has three tiers: (1)~under
two-fold cross-validation by school, the zone model has a higher held-out log-likelihood than the well-mixed
school and than independent classes, with the lower 95\,\% bootstrap bound above zero; (2a)~the freely fitted
shares do not fit better than the measured ones (likelihood ratio at most 5.99), and (2b)~they lie within a
total-variation distance of 0.15 of them; (3)~the observed heterogeneity of attack rates between classes lies
inside the 95\,\% predictive interval of the zone model. The verdict `right structure within tolerance' needs
tiers 1, 2b and 3. On synthetic epidemics a true zone model earned it in 80 to 100\pct{} of replicates, and
a well-mixed school or independent classes in none. The direction of tier 1 was foreseeable: the source
reports that transmission in these schools is dominated by the class.
% src: code/bsc_validation2/a4_zone_level/PREREGISTRATION.md sections 0, 2.3-2.6 (serial-interval kernel: gamma, mean 1.7, sd 1.0, daily bins); Endo et al. 2021, Materials and Methods (Europe PMC PMC8609560, read 1 October 2026: gamma with mean 1.7 and SD 1.0, mean serial interval 2.2 d); Vink et al. 2014 abstract (A(H3N2) 2.2 days); data/v2_numbers.json (zone.power: Z overall 1.00 (A), 0.80 (B); H, X 0.00)

\begin{figure}[tbp]
  \centering
  \includegraphics[width=\textwidth]{v2_zone_en}
  % generated by code/v2_figures.py (zone) from data/bsc_validation2/a4_zone_level/{matsumoto_primary.json, seoul_floors.json} and data/schools/lyon_mixing.json
  \caption{Zone level. (a)~Shares of within-school transmission by level: measured contact durations in a
  primary school in Lyon and shares fitted to the influenza epidemic in Matsumoto. (b)~Held-out
  log-likelihood of the zone model minus that of three competitors, by school. (c)~Heterogeneity of attack
  rates between classes: 95\,\% predictive interval of three models. (d)~Seoul call-centre building: cases
  among the occupants of the other floors predicted from the 11th floor; the bar is a prior-predictive
  interval over an assumed time in the shared lifts and lobby, not a measurement.}
  \label{s8b_evidence:fig:zone}
\end{figure}

\paragraph{Result: the structure, not the numbers.} (Model definitions and power: Section~\ref{appS_prot2:sec:zones}.)
Table~\ref{s8b_evidence:tab:zone} and Fig.~\ref{s8b_evidence:fig:zone} give the outcome, which the
re-check reproduced from data downloaded again, with full-data likelihoods equal to six digits. The zone model
predicts held-out schools far better than a well-mixed school ($+798$ units of log-likelihood, ahead in 27 of
29 schools) and than independent classes ($+211$). The share between classes of the same grade agrees to
two decimals (0.10). Exact agreement of the three shares is rejected (likelihood ratio 42.7): the measured
own-class share is larger than the fitted one, 0.76 against 0.66, and the fitted school-level share is almost
twice the measured one, 0.24 against 0.13. This comparison is confounded. The model has no household term,
whereas the source attributes 41\pct{} of the pupils' infections to households (51\pct{} to school, 8\pct{} to
the community); siblings attend the same school in different grades, so infections between siblings at home
appear as transmission across the school and inflate the fitted school-level share. The shortfall of the
measured school share is therefore probably overstated by this confounding; how much cannot be estimated from
the released table, which has no household identifier, and other unmodelled sources (the external hazard,
contact time as a proxy for transmission) act in directions that are not known. The source itself, with households and the community in its
model, estimated from the same data that transmission to classmates makes up about two-thirds of the
within-school reproduction number when a grade has a single class \citep{endo2021within}; the fitted own-class
share here, 0.66, is of that size. The fitted shares lie within the tolerance (0.103), but not securely: the bootstrap interval
of the distance reaches 0.168, 11\pct{} of resamples exceed the limit, and so does one of the sensitivity
analyses listed in the protocol (no winter break: 0.162). Among the models of the protocol the third tier
excludes only the well-mixed school; the interval of the generic two-level model, simulated afterwards by the
re-check (2.32 to 3.11), does not contain the observed 3.38, a small point in favour of the zone
model. The verdict of the rule, `right structure within tolerance, not exact', is therefore reported as a weak pass.
% src: Endo et al. 2021, Results and Discussion (Europe PMC PMC8609560, read 1 October 2026: school 51.1 %, household 41.3 %, community 7.7 % of student cases; classmates about two-thirds of the within-school reproduction number when a grade has one class); data/v2_numbers.json (zone.posthoc.schools_Z_ahead_of_H [27, 29]; zone.primary_other.LR_F_vs_Z 42.71, w_hat; zone.TV_boot_ci; zone.check.P_TV_gt_0.15 0.107; zone.sensitivity.no_break_switch.TV 0.162); code/bsc_validation2/a4_zone_level/POSTHOC_LOG.md (P4); data/v2_numbers.json (zone.check.disp_C: [2.32, 2.71, 3.11]); notes/VERIFICATION.md (Section 4.9, zone level, items 4, 5, 7; further point 4)

Two comparisons bound what the pass means. First, the generic two-level model with one fitted ratio is not
beaten: the zone model is ahead by 18.8 out of sample, with an interval that includes zero, and the generic
model fits the full data better by 13.6. What the measurement supplies is a ratio that the generic model has to
fit: about 62 implied by the Lyon shares in the Matsumoto class structure, against 37 fitted. Second, a
comparison added during the re-check asks whether measured mixing does better than a guess, in two runs that differ only in the lower
bound of $10^{-9}$ on the external hazard of the protocol. With that bound, a zone model with the round shares
0.7, 0.1 and 0.2 predicts held-out schools better than the measured shares by 19 units (3 to 38), and the
guess two-thirds, one-sixth, one-sixth by 8.9 ($-12.1$ to 34.5, not significant); without it, the guess 0.7, 0.1, 0.2 is better by 34 (11 to 66)
and 0.6, 0.2, 0.2 by 25 ($-17$ to 76). Guesses with 80\pct{} or more in the own class and equal thirds do worse
(all three significantly), and so, in point estimate only, does 0.5, 0.25, 0.25 (by 30, $-30$ to 83;
run without the bound). These guesses were chosen after the fitted shares were known. With that reservation, the measurement is not shown to matter beyond the statement that most
transmission is within the class. Had contact been counted as distinct pairs per day instead of duration, a
choice the protocol excluded in advance, the shares would have been 0.43, 0.14 and 0.43 and the tolerance
would have been missed (0.23). And at the ratio that the data favour, a world in which the generic model is
true earns the full verdict in about 38\pct{} of synthetic replicates (6 of 16), far above the 0 to 7.5\pct{}
obtained when the ratio is drawn at random.
% src: data/v2_numbers.json (zone.D_Z_minus_C; zone.posthoc: AIC, rho_fitted_C.full 36.6, rho_implied_by_Lyon_per_class_quantiles median 61.8; zone.check.hazard_bound."1e-09": D_N702 -19.1 [-37.8, -3.3], D_N2/3 -8.9 [-34.5, 12.1], D_N811 +18.6 (run with the bound); zone.check.D_Z_minus_N_.7_.1_.2: -34.5 [-66.3, -10.6]; zone.check.D_Z_minus_N_.6_.2_.2: -25.3 [-75.8, 17.0]; zone.check.D_Z_minus_N_.5_.25_.25: 29.6 [-29.5, 82.9] (run without the bound); zone.check.TV_naive.Lyon_distinctpairs 0.231); data/checks/zones.json: generic_ratio_model_at_fitted_ratio (6/16 at ratio 37), full_data_log_likelihood_advantage_of_constant_ratio_model (13.6); notes/VERIFICATION.md (Section 4.9, zone level, items 9, 7, 8)

The frequency-dependent normalisation of eq.~\eqref{M-s8b_evidence:eq:zone} fits better than a
density-dependent transfer of the measured contact times, by 37 units of log-likelihood on the full data, in a
sensitivity analysis that the re-check did not re-run. The source had already reported, from the same data,
that the within-school reproduction number hardly depends on class size or the number of classes, that is,
frequency-dependent mixing \citep{endo2021within}; the comparison here agrees with it and adds nothing
independent of it.
% src: data/v2_numbers.json (zone.sensitivity: base_full_data LR_F_vs_Z 42.7, density_dependent_transfer LR_F_vs_Z 116.5); Endo et al. 2021, Results (Europe PMC PMC8609560, read 1 October 2026: R_S "not significantly associated with n or m", "suggesting frequency-dependent mixing"); notes/VERIFICATION.md (Section 4.9, zone level, item 10: unverifiable)

\paragraph{A building by floor: inconclusive.}
The Seoul call-centre building had 94 cases among the 216 people on the 11th floor and 3 among the 927
occupants of the other floors \citep{park2020}. A zone model with one zone per floor and a shared zone of
lifts and lobby, its rate calibrated on the 11th floor alone, predicts a median of 4 cases elsewhere (interval
0--19); a well-mixed building would have had 79. The time spent in the shared zone is not measured, however:
it is a prior of 3 to 30 minutes a day, placed after a rough calculation on the observed counts, and with 30
to 120 minutes the same model would be inconsistent with the data. The check shows that a well-mixed building
is excluded, which needs no model, and nothing more. For the two wings of the 11th floor, and for the
statements of Section~\ref{M-s4_geometry:sec:hotspots} on early growth and on the concentration of risk in
high-efficiency zones, no outbreak with an independent measurement of mixing was found (a cruise ship, a
warship, the meat-processing plant, prisons and nursing homes were considered), and no test was run.
% src: data/v2_numbers.json (zone.seoul: pred_count_median 4, interval [0, 19], well_mixed_expected 78.7, obs_other 3, N_other 927); data/checks/zones.json: seoul_by_floor (30-120 min gives 7 to 81; notes/VERIFICATION.md, Section 4.9, zone level, item 11); code/bsc_validation2/a4_zone_level/PREREGISTRATION.md sections 3, 4

\subsection{All events with a stratum test}

\begin{figure}[tbp]
  \centering
  \includegraphics[width=\textwidth]{v2_summary_scores_en}
  % generated by code/v2_figures.py (summary_scores) from data/v2_numbers.json (first.per_event),
  % data/bsc_validation2/a1_more_outbreaks/10_corrected.json, data/bsc_validation2/a3_tracer_physics/04_evaluation.json and E2_ext_evaluation.json
  \caption{Every event with a stratum test, in both datasets: difference of log scores (a)~between the spatial
  kernel and the well-mixed model and (b)~between the spatial kernel and a constant near/far risk ratio.
  Positive values favour the kernel. Circles: $\Mtwo$ with the mixing coefficient fitted to outbreaks (first
  set: train cohort for vehicles, restaurant for the classroom; second dataset: the six pre-2020 events). In (b)
  the first dataset is compared with the fixed ratios 2 and 3 (digits; introduced after unblinding) and the second
  with the ratio calibrated on the same events as the kernel. Squares: the kernel with its coefficient taken
  from tracer measurements, against the ratios 2 and 3; open squares: its exploratory, outcome-aware application
  to four flights of the second dataset. The restaurant is a calibration event of the first test and has no
  circle; the processing line has no tracer measurement. The kernel is to the right of zero in (a) for all events
  but the coach, and scattered about zero in (b).}
  \label{s8_outbreaks:fig:summary}
\end{figure}
